\documentclass[11pt,a4paper]{article}
\usepackage[T1]{fontenc}
\usepackage[utf8]{inputenc}
\usepackage{lmodern}
\usepackage[polish]{babel}
\usepackage[margin=2.55cm]{geometry}
\usepackage[unicode]{hyperref}
% BEGIN GLOBAL CODE LISTINGS (generated from latex/pl/code-listings.tex)
\usepackage{xcolor}
\usepackage{listings}
\definecolor{codebg}{RGB}{248,249,251}
\definecolor{codeborder}{RGB}{180,191,203}
\definecolor{codecomment}{RGB}{61,113,79}
\definecolor{codekeyword}{RGB}{41,78,132}
\lstdefinelanguage{Julia}{
  morekeywords={abstract,baremodule,begin,break,catch,const,continue,do,else,elseif,
    end,export,false,finally,for,function,global,if,import,in,isa,let,local,macro,
    module,mutable,primitive,quote,return,struct,true,try,type,using,where,while,
    missing,nothing,println,print},
  sensitive=true,morecomment=[l]{\#},morecomment=[n]{\#=}{=\#},
  morestring=[b]",
}
\lstdefinestyle{skrypt}{
  backgroundcolor=\color{codebg},frame=single,rulecolor=\color{codeborder},
  framesep=5pt,frameround=tttt,framerule=.5pt,
  basicstyle=\ttfamily\footnotesize,keywordstyle=\color{codekeyword}\bfseries,
  commentstyle=\color{codecomment}\itshape,stringstyle=\color{purple!65!black},
  numbers=left,numberstyle=\tiny\color{gray!75!black},numbersep=8pt,
  stepnumber=1,breaklines=true,breakatwhitespace=false,
  columns=fullflexible,keepspaces=true,showstringspaces=false,tabsize=4,
  linewidth=.90\linewidth,captionpos=b,aboveskip=11pt,belowskip=12pt,
  xleftmargin=12pt,xrightmargin=2pt,
}
\lstdefinestyle{pythonpde}{style=skrypt,language=Python}
\lstdefinestyle{python}{style=skrypt,language=Python}
\lstdefinestyle{julia}{style=skrypt,language=Julia}
\lstset{style=skrypt}
\renewcommand{\lstlistingname}{Kod}
\renewcommand{\lstlistlistingname}{Spis kodów}
% END GLOBAL CODE LISTINGS

\begin{document}
\section*{Wspólny styl kodu — Python i Julia}
\begin{lstlisting}[language=Python,caption={Obliczanie średniej}]
import numpy as np
# Compute the arithmetic mean.
def mean(values):
    if len(values) == 0:
        raise ValueError("Empty array")
    return np.sum(values) / len(values)
\end{lstlisting}
\begin{lstlisting}[language=Julia,caption={Obliczanie średniej}]
using Statistics
# Compute the arithmetic mean.
function sample_mean(values)
    if isempty(values)
        throw(ArgumentError("Empty array"))
    end
    return mean(values)
end
#= A block comment
   #= nested =#
   across multiple lines. =#
println(sample_mean([1, 2, 3, 4]))
\end{lstlisting}
\begin{lstlisting}[style=pythonpde]
# Compatibility with earlier PDE chapters.
print("Python")
\end{lstlisting}
\end{document}
